% ========== CHAPTER 6 COVER: MICROKERNEL ARCHITECTURE ==========
% Creative and elegant chapter introduction with sophisticated visual elements
% Following the established 4-element pattern
% Requirements: 7.1, 7.2, 7.3, 7.4, 7.5

\vspace{2cm}

% 1. BIG CHAPTER NUMBER in content-reflective shape (gear pattern for system architecture)
\begin{center}
\begin{tikzpicture}
    % Gear/mechanical pattern representing system architecture
    \fill[brandPrimary!15] (0,0) circle (2.2);
    % Inner gear teeth pattern
    \foreach \angle in {0,30,60,90,120,150,180,210,240,270,300,330} {
        \fill[brandPrimary!20] (\angle:1.8) -- (\angle+15:2.2) -- (\angle+30:1.8) -- cycle;
    }
    % Central hub
    \fill[brandPrimary!25] (0,0) circle (1.2);
    % Chapter number
    \node[font=\fontsize{80}{80}\selectfont\bfseries, color=brandPrimary] at (0,0) {6};
\end{tikzpicture}
\end{center}

\vspace{1cm}

% 2. CHAPTER TOPIC with elegant typography
\begin{center}
{\fontsize{24}{28}\selectfont\textcolor{brandPrimary}{\textit{Microkernel}}} \\
\vspace{0.2cm}
{\fontsize{28}{32}\selectfont\textcolor{brandSecondary}{\textbf{Architecture}}}
\end{center}

\vspace{0.5cm}

% 3. NICE SEPARATOR with architectural symbols
\begin{center}
\textcolor{brandAccent}{
\raisebox{-0.2ex}{\Large$\hexagon$} \quad
\rule{1.8cm}{0.8pt} \quad
\raisebox{-0.2ex}{\Large$\star$} \quad
\rule{3.5cm}{0.8pt} \quad
\raisebox{-0.2ex}{\Large$\hexagon$} \quad
\rule{1.8cm}{0.8pt} \quad
\raisebox{-0.2ex}{\Large$\star$}
}
\end{center}

\vspace{1.5cm}

% 4. "THIS CHAPTER EXPLORES" content box
\begin{tcolorbox}[
    enhanced,
    colback=white,
    colframe=brandPrimary!40,
    boxrule=1.5pt,
    arc=15pt,
    drop shadow={brandPrimary!25},
    left=1.5cm,
    right=1.5cm,
    top=1.2cm,
    bottom=1.2cm
]
\begin{center}
{\Large\textbf{\textcolor{brandPrimary}{This chapter explores}}}
\end{center}

\vspace{0.8cm}

\begin{itemize}[leftmargin=1.2cm, itemsep=0.4cm, label={\textcolor{brandAccent}{$\triangleright$}}]
    \item {\large Microkernel design philosophy and the "Minimal Core, Maximum Potential" principle}
    \item {\large Extension execution models: The Pit's in-process performance vs The Grand Stage's isolation}
    \item {\large Kernel bootstrap process with dependency resolution and failure handling mechanisms}
    \item {\large SPFR (Symphony Playbook Failure Recovery) framework for system reliability}
    \item {\large Performance characteristics and trade-offs between safety and speed}
    \item {\large Circuit breakers, retry logic, and health monitoring for resilient orchestration}
\end{itemize}
\end{tcolorbox}

\clearpage